\documentclass[11pt,a4paper]{article}
% =====================================================================
%  Shared preamble for the Part V lecture notes (lectures 19-22).
%
%  These four lectures sit outside Geron, so the notes ARE the primary
%  source and are examined on the same terms as the textbook parts.
%  Everything here is chosen to make that readable on paper: one column,
%  generous measure, and the same six-colour palette as the slides so a
%  student moving between the two is not re-learning what red means.
%
%  Build with pdflatex (twice, for the toc and the refs):
%      cd notes && latexmk -pdf lecture-19.tex
%  or  make -C notes
% =====================================================================
\usepackage[T1]{fontenc}
\usepackage[utf8]{inputenc}
\usepackage{newtxtext,newtxmath}      % Times-like: prints well, reads well
\usepackage[scaled=0.86]{inconsolata} % code, at a size that sits beside Times
\usepackage{microtype}
% newtx defines \Bbbk, \openbox, \proof and \endproof; amssymb and amsthm
% then redefine them and abort. Freeing the four names before those two
% packages load is the standard fix, and it costs nothing here: we use
% amsthm only for \newtheorem.
\let\Bbbk\relax
\usepackage{amsmath,amssymb}
\let\openbox\relax
\let\proof\relax
\let\endproof\relax
\usepackage{amsthm}
\usepackage{booktabs}
\usepackage{enumitem}
\usepackage{xcolor}
\usepackage{tcolorbox}
\usepackage{titlesec}
\usepackage{graphicx}
\usepackage[hidelinks]{hyperref}
\tcbuselibrary{skins,breakable}

% ---- the course palette, from assets/css/custom.css -------------------
\definecolor{aimlprimary}{HTML}{0B3D62}   % deep blue  - structure
\definecolor{aimlaccent} {HTML}{C0392B}   % brick red  - failures
\definecolor{aimlsuccess}{HTML}{14663A}   % green      - repairs
\definecolor{aimlmath}   {HTML}{6C3483}   % purple     - the derivation
\definecolor{aimlmuted}  {HTML}{4B5563}
\definecolor{aimlrule}   {HTML}{B0BCC7}
\definecolor{aimlsoft}   {HTML}{F4F7F9}

\usepackage[a4paper,margin=32mm,top=30mm,bottom=30mm]{geometry}
\setlength{\parskip}{0.45\baselineskip}
\setlength{\parindent}{0pt}
\linespread{1.06}

% ---- headings --------------------------------------------------------
\titleformat{\section}{\Large\bfseries\color{aimlprimary}}{\thesection}{0.7em}{}
\titleformat{\subsection}{\large\bfseries\color{aimlprimary}}{\thesubsection}{0.7em}{}
\titleformat{\subsubsection}{\bfseries\color{aimlmuted}}{}{0em}{}
\titlespacing*{\section}{0pt}{1.6\baselineskip}{0.5\baselineskip}
\titlespacing*{\subsection}{0pt}{1.2\baselineskip}{0.35\baselineskip}

% ---- boxes -----------------------------------------------------------
% One box per role, and the role is the colour. A student who has seen the
% slides already knows what each one means before reading a word of it.
\newtcolorbox{keybox}[1][]{
  breakable, enhanced, colback=aimlsoft, colframe=aimlprimary,
  boxrule=0.4pt, left=8pt, right=8pt, top=6pt, bottom=6pt,
  fonttitle=\bfseries\small, coltitle=white, colbacktitle=aimlprimary,
  attach boxed title to top left={yshift=-2mm, xshift=4mm},
  boxed title style={boxrule=0pt, arc=1pt}, #1}

\newtcolorbox{failbox}[1][]{
  breakable, enhanced, colback=white, colframe=aimlaccent,
  boxrule=0.9pt, left=8pt, right=8pt, top=6pt, bottom=6pt,
  fonttitle=\bfseries\small, coltitle=white, colbacktitle=aimlaccent,
  attach boxed title to top left={yshift=-2mm, xshift=4mm},
  boxed title style={boxrule=0pt, arc=1pt}, #1}

\newtcolorbox{derivbox}[1][]{
  breakable, enhanced, colback=white, colframe=aimlmath,
  boxrule=0.6pt, left=8pt, right=8pt, top=6pt, bottom=6pt,
  fonttitle=\bfseries\small, coltitle=white, colbacktitle=aimlmath,
  attach boxed title to top left={yshift=-2mm, xshift=4mm},
  boxed title style={boxrule=0pt, arc=1pt}, #1}

% An exercise. Deliberately the same shape as the notebook's `try` field:
% one change, and what should happen to the output.
\newtcolorbox{trybox}[1][]{
  breakable, enhanced, colback=white, colframe=aimlsuccess,
  boxrule=0.5pt, left=8pt, right=8pt, top=6pt, bottom=6pt,
  fonttitle=\bfseries\small, coltitle=white, colbacktitle=aimlsuccess,
  attach boxed title to top left={yshift=-2mm, xshift=4mm},
  boxed title style={boxrule=0pt, arc=1pt}, #1}

% ---- small semantics -------------------------------------------------
\newcommand{\scope}[1]{\textcolor{aimlmuted}{\small #1}}
\newcommand{\term}[1]{\textbf{#1}}
\newcommand{\code}[1]{\texttt{\small #1}}
\newcommand{\lecture}[1]{Lecture~#1}

\newtheorem{definition}{Definition}[section]
\newtheorem{proposition}[definition]{Proposition}

% ---- title block -----------------------------------------------------
% \notestitle{19}{Information retrieval: the lexical foundation}{SciFact (BEIR)}
\newcommand{\notestitle}[3]{%
  \thispagestyle{empty}
  \begin{flushleft}
    {\color{aimlmuted}\small\sffamily
     APPLICATIONS OF MACHINE LEARNING \quad$\cdot$\quad
     BSc Mathematics of Artificial Intelligence}\\[2mm]
    {\color{aimlrule}\rule{\textwidth}{0.8pt}}\\[4mm]
    {\color{aimlmuted}\large Lecture #1 \quad$\cdot$\quad Part V \quad$\cdot$\quad
     Extended notes}\\[3mm]
    {\Huge\bfseries\color{aimlprimary}\baselineskip=1.15\baselineskip #2\par}\vspace{4mm}
    {\color{aimlmuted}\large Dataset: #3}\\[3mm]
    {\color{aimlrule}\rule{\textwidth}{0.8pt}}
  \end{flushleft}
  \vspace{2mm}
  \begin{keybox}[title={Why these notes exist}]
    Lectures 19--22 sit outside G\'eron, the course's primary text. For those
    four lectures \emph{these notes are the primary source}, and they are
    examined on exactly the same terms as the textbook parts. Everything in
    the slide deck for this lecture is here, with the arguments written out at
    length rather than compressed onto a slide; the numbers are the ones the
    accompanying notebook prints, and every one of them is a measurement on
    the corpus named above rather than a figure quoted from a paper.
  \end{keybox}
  \vspace{1mm}
  {\small\tableofcontents}
  \vspace{2mm}
  \noindent{\color{aimlrule}\rule{\textwidth}{0.4pt}}
  \vspace{1mm}
}


\begin{document}
\notestitlefor{16}{IV}{Recurrent networks}{Chicago transit ridership}{Chapter 13}

\section{What a linear model on lags cannot do}

Lecture 15 established the ground rules and a number to beat: evaluation is
forward in time, never a random split; the seasonal-naive forecast is genuinely
hard to beat and the autocorrelation says why; and a linear model on lagged
features beats it, modestly.

That model has a \emph{fixed window}. Anything outside the last 56 days is
invisible to it, and every lag gets one weight regardless of context. Take 365
lags and it has 365 weights, each fitted from the same amount of data, each
learning one number about one position --- and the weight on lag 7 learns
nothing from the weight on lag 14, though both are ``a week ago''.

A recurrent network keeps a \emph{state} instead, and that state is not bounded
by a window:
\[ \mathbf{h}_t = \phi\!\left(\mathbf{W}_x^{\mathsf T}\mathbf{x}_t
   + \mathbf{W}_h^{\mathsf T}\mathbf{h}_{t-1} + \mathbf{b}\right). \]
This is Lecture 9's $\phi(\mathbf{X}\mathbf{W} + \mathbf{b})$ with one extra
term, and that term is the whole idea. $\mathbf{W}_x$ and $\mathbf{W}_h$ do not
depend on $t$: one set of weights applied at every step, so reach is decoupled
from parameter count --- the same argument as weight sharing in Lecture 12,
along time instead of across an image.

\begin{keybox}[title={The state is a summary, and that is a budget}]
$\mathbf{h}_t$ is everything the network has decided to remember about
$x_1,\dots,x_t$, at \emph{fixed width} whatever the sequence length --- so it
is lossy by construction. What it keeps is learned, not designed: nothing tells
the cell that ``the same weekday last week'' matters.

Two sequences with the same state are indistinguishable to everything
downstream, which is the definition of a sufficient statistic, and the network
is trying to learn one. It is also the honest limit: if the task needs more
bits than the state has, no amount of training supplies them.
\end{keybox}

\subsection{Unrolling, and the problem it creates}

Draw the same cell once per time step and you get a network 56 layers deep with
one set of weights, reused. Then differentiate it exactly as Lecture 10 taught
--- \term{backpropagation through time} needs nothing new.

\begin{failbox}[title={Failure condition --- one matrix raised to the length of the sequence}]
It is multiplied by the same factor 56 times. Here the factor is \emph{the same
matrix} every step, because the weights are shared --- so the gradient is
governed by a power of one matrix, and powers of a matrix either vanish or
explode unless the spectral radius sits almost exactly at one.

Initialisation cannot fix this the way it did in Lecture 11, where the layers
were independent. Something has to change in the cell itself.
\end{failbox}

That is also why $\tanh$ rather than ReLU: a bounded activation keeps the
repeated product from exploding.

The model is eleven lines --- an \code{nn.RNN} and a linear head applied to
\code{outputs[:, -1]}, the hidden state after the final step. The head is not
optional: the state has 32 components and the target has one, and $\tanh$
outputs lie in $(-1,1)$ while our scaled targets sometimes exceed that.
\code{batch\_first=True} because our batches are [batch, time, series];
without it the module silently reads the first dimension as time.

The training loop is Lecture 10's five lines, unchanged. The loss is Huber ---
quadratic near zero, linear far away --- because we are optimising a proxy for
MAE, and minimising the absolute error directly gives a gradient that never
gets smaller as you approach the optimum. \emph{The metric and the training
loss are allowed to differ}; Lecture 2 said so, and here is the case.

\begin{keybox}[title={Shuffle, but shuffle the right thing}]
\code{DataLoader(train\_set, shuffle=True)} shuffles the \emph{windows}. It
does not shuffle the values inside a window --- the sequence within each row is
untouched.

Consecutive windows share almost all their content: one starting on day $t$ and
one on day $t+1$ differ in two values out of fifty-six. Without shuffling, a
mini-batch is fifty near-identical examples, its gradient is one example's
gradient with less noise, and training stalls.

And do not let that shuffling cross the split. Windows are shuffled within the
training period only. Shuffling and leakage are different operations, and it is
easy to conflate them.
\end{keybox}

On identical shuffled folds the recurrent network scores 38{,}224, against
55{,}385 for copying last week and 44{,}925 for the linear model --- 31.0\%
better than the baseline and 15\% better than the linear model. Both beat the
baseline by more than the fold-to-fold spread, so the ordering is real, and the
recurrent model wins because it can use the \emph{order} of the window rather
than a weighted sum of it.

\begin{failbox}[title={Failure condition --- and the assumption, stated rather than left implied}]
Those folds are shuffled, so \emph{none of those numbers is a forecast}.
Lecture 15 measured what that is worth. They are comparable with each other
because the folds are identical, and with nothing else.

Every forward number in the rest of this lecture is on the honest ladder.
\end{failbox}

\section{Four improvements, measured}

\subsection{Deeper --- which failed}

One argument, \code{num\_layers=3}. The simple RNN forecasts at 49{,}073 and
the three-layer one at 53{,}491.

It got worse, and not from overfitting: its error on the \emph{training} rows
is higher too, 40{,}630 against 35{,}972. Three times the parameters fitted on
the same 1{,}040 windows, and gradients now traversing 56 time steps
\emph{and} three layers --- Lecture 11's vanishing signal compounded.

\emph{Report the negative result. A ladder of improvements in which every rung
works is a ladder somebody edited.}

\subsection{More series --- which worked}

We have been forecasting rail from rail alone while a bus series sits unused,
and a \code{day\_type} column we established is known in advance. Adding both
changes the architecture by one number, \code{input\_size=5}.

\begin{failbox}[title={Wait --- that is a \code{shift(-1)}}]
Lecture 15 spent ten minutes on a random split that let a model train on rows
surrounding what it was scored on. This one reads the future \emph{deliberately}.

The test that separates them: \emph{is the value available at prediction time?}
The type of the day we are forecasting is on the public calendar and has been
for a year. Its ridership is not.

This is the one question about leakage that no static check can answer for you.
It is not a property of the code; it is a property of the world.
\end{failbox}

\begin{center}
\small
\begin{tabular}{lrl}
\toprule
\textbf{Model} & \textbf{Forecast MAE} & \textbf{vs copying last week} \\
\midrule
Copy last week & 64{,}815 & --- \\
Simple RNN, rail only & 49{,}073 & 24\% better \\
Simple RNN, five inputs & 40{,}046 & 38\% better \\
\bottomrule
\end{tabular}
\end{center}

Eighteen per cent off the error for one line of feature engineering and one
changed constructor argument. Most of it is the day type: the network no longer
has to infer that a holiday is coming from the shape of the previous eight
weeks --- which is impossible --- because we told it.

\subsection{Further ahead}

Three strategies, all defensible, all to be evaluated forward in time at the
horizon you actually need. \term{Recursive} predicts one step and feeds it
back, so errors compound and the model never trained on its own output.
\term{Direct} fits a separate model per horizon, so nothing is shared.
\term{Sequence-to-sequence} emits $k$ outputs at once.

A fourth variant predicts fourteen at \emph{every} step, taking the head over
all outputs rather than the last. That is not cheating: a recurrent cell at
step $t$ has seen steps 1 to $t$ and nothing more --- it is \term{causal}. The
loss simply gains a term at every step, so many more gradients flow and they
travel a shorter distance. After training, only the last step's output is used.

Measured, the error grows with distance from 35{,}997 at one day to 46{,}586 at
fourteen, and stays below the no-model baseline throughout.

\begin{keybox}[title={Teacher forcing, and what it costs}]
Feeding the decoder the true previous token trains fast and stable, but the
model has never seen its own mistakes; feeding its own prediction matches
deployment and makes early training noise.

\term{Exposure bias} is the mismatch: a model trained only on true prefixes
meets a distribution at prediction time it never saw --- its own errors,
compounding --- which is precisely the recursive strategy's weakness.
\end{keybox}

\begin{failbox}[title={Failure condition --- bidirectional layers, and the one question}]
Running a second layer backwards lets every position see the whole sequence. A
backward pass \emph{reads the future}: fine when the whole sequence is
available at prediction time, and leakage when it is not.

A bidirectional forecaster is Lecture 15's random split wearing a different
hat --- the number goes up and the system cannot be deployed. Ask: at
prediction time, do I have the rest of the sequence? For text classification,
usually yes. For forecasting, never.
\end{failbox}

\subsection{Gates}

A plain cell multiplies the state at every step, so
$\partial\mathbf{h}_t/\partial\mathbf{h}_{t-1} =
\operatorname{diag}(\phi')\mathbf{W}_h^{\mathsf T}$ and over $T$ steps the
gradient carries $T$ copies of one matrix.

\begin{derivbox}[title={What a gated cell does instead}]
An LSTM keeps a separate cell state updated \emph{additively},
\[ \mathbf{c}_t = \mathbf{f}_t \odot \mathbf{c}_{t-1}
                + \mathbf{i}_t \odot \tilde{\mathbf{c}}_t. \]
When the forget gate is near one,
$\partial\mathbf{c}_t/\partial\mathbf{c}_{t-1} \approx \mathbf{1}$ and the
gradient passes through almost unchanged.

That is the entire trick: \emph{a path along which the gradient is added to
rather than multiplied}.
\end{derivbox}

\begin{center}
\small
\begin{tabular}{lll}
\toprule
\textbf{Gate} & \textbf{Form} & \textbf{What it decides} \\
\midrule
forget $\mathbf{f}_t$ & sigmoid & how much of the old cell state to keep \\
input $\mathbf{i}_t$ & sigmoid & how much of the candidate to write \\
candidate $\tilde{\mathbf{c}}_t$ & $\tanh$ & what to write, if anything \\
output $\mathbf{o}_t$ & sigmoid & how much of the cell state to expose \\
\bottomrule
\end{tabular}
\end{center}

Every gate is a sigmoid of the same two inputs with its own weights, so a
sigmoid in $[0,1]$ is a \emph{soft switch} and the network learns when to open
it. A GRU merges forget and input into one update gate --- what is kept is what
is not written --- and has no separate cell state, keeping the additive
near-identity path with about three quarters of the parameters.

\begin{center}
\small
\begin{tabular}{lrr}
\toprule
\textbf{Model, all five inputs} & \textbf{Forecast MAE} &
\textbf{Training MAE} \\
\midrule
Simple RNN & 40{,}046 & 27{,}035 \\
LSTM & 37{,}542 & 24{,}540 \\
GRU & 34{,}639 & 20{,}356 \\
\bottomrule
\end{tabular}
\end{center}

The GRU wins, and the second column says why to be careful: it fits the
training rows 25\% better than the plain cell but the forecast only 14\%
better. That gap is overfitting, and it is the next thing to attack.

A one-dimensional convolution with stride 2 before the recurrent layer lets the
same number of steps cover twice the history --- and on this series it does not
pay, 39{,}792 against 35{,}997 at one day and 55{,}148 against 46{,}586 at
fourteen. Reported because we measured it, not because it won.

\begin{keybox}[title={Gradient clipping is not optional here}]
The gradient through $T$ steps carries a power of one matrix. When its spectral
radius exceeds one even slightly the norm does not drift upwards --- it
explodes over a handful of steps, and a single bad batch can move every weight
somewhere training never recovers from.

Clip by \emph{global norm}, not per parameter: the direction of the update is
information and only its length is the problem.
\end{keybox}

\section{The whole application}

\begin{center}
\small
\begin{tabular}{lrr}
\toprule
\textbf{Forecast} & \textbf{What we reported} & \textbf{What it does} \\
\midrule
Copy last week & 55{,}385 & 64{,}815 \\
Linear, 56 lags & 44{,}925 & 56{,}446 \\
Simple RNN & 38{,}224 & 49{,}073 \\
GRU, five inputs & --- & 34{,}639 \\
\bottomrule
\end{tabular}
\end{center}

The two middle rows are the same models measured twice. The bottom row is the
model we would deploy, and it was never measured the wrong way. \emph{The final
number is better than the one we committed to --- but we reached it by first
admitting that the committed one was fiction.}

Is 34{,}639 good? It is 47\% better than a system a transit authority could run
with no data scientist at all, and its typical miss is around 10.4\% of the day
it forecasts, against 16.5\% for copying last week \emph{on those same 151
days}. Lecture 15's 12.0\% was the naive forecast over all 1{,}191 pool days,
which is a different set of days and not the comparison to make here.

\begin{keybox}[title={What a defensible report says}]
``Forecasting the five months after the training cut, and never seeing them,
the model's mean absolute error is 34{,}639 boardings a day against 64{,}815
for copying the same day of the previous week. Both figures are on identical
rows. The gain is concentrated on holidays.''

The protocol, the baseline, the rows, and where the gain comes from --- four
things the number alone does not carry.
\end{keybox}

\subsection{Red-team it: what happens when the world changes}

Fit on everything to the end of 2019 and forecast the four months from March
2020.

\begin{center}
\small
\begin{tabular}{lr}
\toprule
& \textbf{MAE} \\
\midrule
The linear model, trained to the end of 2019 & 84{,}380 \\
Copying last week, on the same 122 days & 36{,}282 \\
The actual mean level over that period & 141{,}404 \\
\bottomrule
\end{tabular}
\end{center}

The model is 2.3 times worse than doing nothing, and its error is six tenths of
the entire level.

\begin{failbox}[title={Failure condition --- and the shuffled split cannot see it at all}]
The same shuffled five-fold cross-validation over 2016--2021, spanning the
collapse, reports 36{,}554 --- \emph{better} than before, because the post-2020
days are easy to interpolate: flat, low, and surrounded by other flat low days
that are in the training set.

A random split does not merely overstate performance. It is \emph{structurally
incapable} of detecting a regime change, which is the failure that actually
ends forecasting systems.
\end{failbox}

To be fair to the model: March 2020 is not a leak and not its fault. No
forecaster could have produced it from eight weeks of history, and copying last
week wins only because it adapts within a week of the collapse. The lesson is
not that the model is bad --- it is that \emph{an honest protocol told you the
model had failed, on the day it failed, and the dishonest one reported an
improvement}. Which is why a deployed forecaster is monitored against its own
naive baseline, continuously: if the margin disappears, something has changed
in the world.

\begin{keybox}[title={Temporal leakage checklist}]
\begin{enumerate}[leftmargin=1.6em,itemsep=1pt]
  \item Is the split by \emph{date}, not by row?
  \item Does every feature exist at the moment the forecast is made? Name the
        moment.
  \item Were any statistics computed over the whole series before the cut?
  \item Do the training windows overlap the scored windows? Purge them.
  \item Is the baseline computed on identical rows to the model?
  \item Does the reported score come from more than one origin?
\end{enumerate}
The third is Lecture 2's leak meeting today's.
\end{keybox}

\section{What recurrence still cannot do}

It is \emph{sequential}: step $t$ needs step $t-1$, so the computation cannot
be parallelised along the sequence. Batching parallelises across sequences,
never along one, so training time grows linearly with sequence length on any
hardware. On a 56-step window nobody cares; on a document of ten thousand
tokens it is the binding constraint, and it is the practical reason --- more
than the gradient --- that attention replaced recurrence.

The state is a \emph{bottleneck}: everything remembered must fit in one
fixed-width vector. And long-range dependence is still hard --- gates make the
gradient survive, but they do not make position 400 as accessible as position
2.

Let every position look directly at every other, with learned weights, and all
three problems go at once: the computation parallelises, there is no single
state to squeeze through, and position 400 is exactly as far away as position
2. That is attention, and it is Lecture 18.

\section{The notebook}

\begin{trybox}[title={What to change}]
Lengthen the window from 56 steps to 200 and re-run the plain RNN. It gets
\emph{worse}, not better, and the gradient probe --- Lecture 11's, along time
instead of depth --- says why.
\end{trybox}

\section{Exercises}

\begin{enumerate}[leftmargin=1.6em,itemsep=4pt]
  \item Explain why unrolling a recurrent cell over $T$ steps makes the
        gradient harder to control than a $T$-layer feed-forward network, and
        why initialisation cannot repair it. \hfill \scope{[5]}
  \item State the update of an LSTM's cell state, and explain in one sentence
        why its additive form is what lets the gradient survive.
        \hfill \scope{[4]}
  \item A three-layer recurrent network scores worse than a one-layer one on
        both the forecast \emph{and} the training rows. State what that rules
        out, and what it suggests instead. \hfill \scope{[4]}
  \item A feature is built with \code{shift(-1)}. Give the single question that
        decides whether this is leakage, and explain why no static check can
        answer it. \hfill \scope{[4]}
  \item A shuffled cross-validation over a period spanning a regime change
        reports a better score than one that excludes it. Explain the
        mechanism, and what it implies about monitoring a deployed forecaster.
        \hfill \scope{[5]}
\end{enumerate}

\section*{Summary}
\addcontentsline{toc}{section}{Summary}

A recurrent cell is Lecture 9's layer with one extra term, and that term
decouples reach from parameter count --- weight sharing along time instead of
across an image. The state is a fixed-width lossy summary, which is a budget:
if the task needs more bits than it holds, no amount of training supplies them.
Unrolling makes training ordinary backpropagation on a 56-layer network with
\emph{shared} weights, and that sharing is the problem, because the gradient
then carries a power of one matrix rather than a product of independent ones.

Four improvements, measured forward in time, and one of them failed: a
three-layer network scored worse on the forecast \emph{and} on the training
rows, which rules out overfitting and points at optimisation. Adding the bus
series and the next day's type --- a \code{shift(-1)} that is legitimate
because the calendar is public --- took the error down 18\%. Gates buy an
additive path along which the gradient is added to rather than multiplied, and
the GRU reached 34{,}639.

The ladder ends 47\% better than copying last week, and the honest ladder is
not the one we first reported: the two middle rows are the same models measured
twice. Then the red team --- a model trained to the end of 2019 is 2.3 times
worse than doing nothing on the months from March 2020, while the shuffled
protocol reports an \emph{improvement} over that same period, because flat low
days surrounded by flat low days are easy to interpolate. A random split does
not merely overstate performance; it cannot see a regime change at all.

\end{document}
